\documentclass[aspectratio=169,11pt]{beamer}

% Packages
\usepackage[utf8]{inputenc}
\usepackage[T1]{fontenc}
\usepackage{lmodern}
\usepackage{microtype}
\usepackage{amsmath,amssymb}
\usepackage{booktabs}
\usepackage{graphicx}
\usepackage{tikz}
\usepackage{pgfplots}
\usepackage{xcolor}
\usepackage{ragged2e}
\usepackage{colortbl}
\usepackage{array}
\usepackage{hyperref}
\usepackage{adjustbox}

\usetikzlibrary{shapes.geometric, arrows.meta, positioning, calc, backgrounds, decorations.pathreplacing, shadows, fadings}
\pgfplotsset{compat=1.18}

% --- Color Palette: Warm Professional ---
\definecolor{DeepNavy}{HTML}{2E4057}
\definecolor{Teal}{HTML}{048A81}
\definecolor{WarmOrange}{HTML}{E85D04}
\definecolor{SoftPurple}{HTML}{9D4EDD}
\definecolor{WarmGray}{HTML}{6C757D}
\definecolor{LightGray}{HTML}{E9ECEF}
\definecolor{Cream}{HTML}{FBF8F1}
\definecolor{DeepRed}{HTML}{D62828}
\definecolor{Gold}{HTML}{D4A03A}
\definecolor{SoftWhite}{HTML}{FAFAFA}

% Beamer colors
\setbeamercolor{normal text}{fg=DeepNavy, bg=SoftWhite}
\setbeamercolor{structure}{fg=DeepNavy}
\setbeamercolor{alerted text}{fg=DeepRed}
\setbeamercolor{example text}{fg=Teal}
\setbeamercolor{frametitle}{fg=DeepNavy, bg=SoftWhite}
\setbeamercolor{title}{fg=DeepNavy}
\setbeamercolor{subtitle}{fg=WarmGray}
\setbeamercolor{block title}{bg=DeepNavy, fg=SoftWhite}
\setbeamercolor{block body}{bg=LightGray, fg=DeepNavy}
\setbeamercolor{itemize item}{fg=WarmOrange}
\setbeamercolor{itemize subitem}{fg=Gold}

% Fonts
\usefonttheme{professionalfonts}
\setbeamerfont{title}{size=\huge, series=\bfseries}
\setbeamerfont{frametitle}{size=\Large, series=\bfseries}

% Strip chrome
\setbeamertemplate{navigation symbols}{}
\setbeamertemplate{headline}{}

% Minimal footline
\setbeamertemplate{footline}{%
    \hfill
    \begin{beamercolorbox}[wd=3cm,ht=2.5ex,dp=1ex,right,rightskip=0.6cm]{page number}
        \usebeamercolor[fg]{WarmGray}\scriptsize\insertframenumber
    \end{beamercolorbox}
    \vspace{0.3cm}
}

% Bullets
\setbeamertemplate{itemize item}{\footnotesize\raisebox{0.3ex}{\tikz\fill[WarmOrange] (0,0) circle (0.45ex);}}
\setbeamertemplate{itemize subitem}{\footnotesize\raisebox{0.3ex}{\tikz\fill[Gold] (0,0) circle (0.35ex);}}

% Custom commands
\newcommand{\transitionslide}[2]{%
    {\setbeamercolor{normal text}{fg=SoftWhite,bg=DeepNavy}
    \begin{frame}[plain]\vfill\begin{center}
        {\huge\bfseries\textcolor{Gold}{#1}}\\[0.6cm]
        {\large\textcolor{LightGray}{#2}}
    \end{center}\vfill\end{frame}}
}

\begin{document}

\transitionslide{Technical Supplement}{Joint Estimation via GMM}

% --- SLIDE 1 ---
\begin{frame}
    \frametitle{The Identification Problem: Circularity}
    
    \vspace{0.3cm}
    Recall the structural equations:
    \begin{align*}
        w_t &= \psi L_{St} + \varepsilon_t \tag{Supply} \\
        L_{it} &= \phi^d w_t + \lambda_i \eta_t + u_{it} \tag{Demand}
    \end{align*}
    
    \vspace{0.3cm}
    \begin{itemize}
        \item \textbf{The Challenge with PCA:} We cannot perfectly extract $\lambda_i$ via PCA on $L_{it}$ without first knowing the wage effect $\phi^d w_t$. 
        \item \textbf{The Circularity:} We cannot know $\phi^d$ without a valid instrument $z_t$. But constructing $z_t$ requires knowing $\lambda_i$.
        \item \textbf{The Solution:} We must estimate the structural elasticities ($\psi, \phi^d$) and the latent factor loadings ($\lambda_i$) \textbf{jointly}.
    \end{itemize}
\end{frame}

% --- SLIDE 2 ---
\begin{frame}
    \frametitle{The Core Moment Conditions}
    
    \vspace{0.5cm}
    GMM identifies parameters by matching theoretical moments to sample moments. The core idea is that a valid instrument $z_t(\lambda)$ must be orthogonal to aggregate shocks.
    
    \vspace{0.8cm}
    \textbf{1. Supply Moment (Identifying $\psi$):} \\
    The instrument must be uncorrelated with the supply shock $\varepsilon_t$.
    \begin{equation*}
        \mathbb{E}[(w_t - \psi L_{St}) z_t(\lambda)] = 0
    \end{equation*}
    
    \vspace{0.5cm}
    \textbf{2. Demand Moment (Identifying $\phi^d$):} \\
    The instrument must be uncorrelated with the demand shock $\eta_t$. Using the equal-weighted average demand:
    \begin{equation*}
        \mathbb{E}[(L_{Et} - \phi^d w_t) z_t(\lambda)] = 0
    \end{equation*}
\end{frame}

% --- SLIDE 3 ---
\begin{frame}
    \frametitle{The GMM Objective Function}
    
    \vspace{0.5cm}
    Let $\theta = (\psi, \phi^d)$ be the structural parameters, and let the sample moments be:
    \begin{align*}
        g_{1,T}(\theta, \lambda) &= \frac{1}{T} \sum (w_t - \psi L_{St}) z_t(\lambda) \\
        g_{2,T}(\theta, \lambda) &= \frac{1}{T} \sum (L_{Et} - \phi^d w_t) z_t(\lambda)
    \end{align*}
    
    \vspace{0.5cm}
    Let $g_T(\theta, \lambda) = [g_{1,T}, g_{2,T}]'$. The GMM estimator seeks $(\theta, \lambda)$ to minimize:
    \begin{equation*}
        J(\theta, \lambda) = g_T(\theta, \lambda)' W g_T(\theta, \lambda)
    \end{equation*}
    where $W$ is a positive definite weighting matrix.
\end{frame}

% --- SLIDE 4 ---
\begin{frame}[shrink=10]
    \frametitle{Concentrated GMM: The Iterative Loop (Part 1)}
    
    \vspace{0.3cm}
    Simultaneous optimization over elasticities and $N$ loadings is difficult. Gabaix and Koijen (2022) use \textbf{Concentrated GMM} (Iterative IV).
    
    \vspace{0.5cm}
    \begin{enumerate}
        \item \textbf{Initialize:} Start with an initial guess for the elasticities, $\theta^{(0)} = (\psi^{(0)}, \phi^{d(0)})$ (e.g., from OLS).
        
        \vspace{0.3cm}
        \item \textbf{Profile out Endogeneity:} Given $\phi^{d(k)}$, compute partialled-out labor demand:
        \begin{equation*}
            \tilde{L}_{it}^{(k)} = L_{it} - \phi^{d(k)} w_t = \lambda_i \eta_t + u_{it}
        \end{equation*}
        
        \vspace{0.3cm}
        \item \textbf{Extract Loadings via PCA:} The matrix $\tilde{L}^{(k)}$ now has a pure factor structure. Extract principal eigenvector to get updated loadings $\lambda^{(k+1)}$.
    \end{enumerate}
\end{frame}

% --- SLIDE 5 ---
\begin{frame}
    \frametitle{Concentrated GMM: The Iterative Loop (Part 2)}
    
    \vspace{0.5cm}
    \begin{enumerate}
        \setcounter{enumi}{3}
        \item \textbf{Construct Instrument:} Using $\lambda^{(k+1)}$, compute the projection matrix $Q$, the optimal weights $\Gamma^{(k+1)} = S'Q$, and the instrument:
        \begin{equation*}
            z_t^{(k+1)} = \sum \Gamma_i^{(k+1)} L_{it}
        \end{equation*}
        
        \vspace{0.3cm}
        \item \textbf{Update Parameters (IV):} Update elasticities by setting sample moments to zero (equivalent to 2SLS using the new instrument):
        \begin{align*}
            \psi^{(k+1)} = \frac{\sum w_t z_t^{(k+1)}}{\sum L_{St} z_t^{(k+1)}}, \quad \phi^{d(k+1)} = \frac{\sum L_{Et} z_t^{(k+1)}}{\sum w_t z_t^{(k+1)}}
        \end{align*}
        
        \vspace{0.3cm}
        \item \textbf{Iterate:} Repeat steps 2--5 until parameters converge.
    \end{enumerate}
\end{frame}

% --- SLIDE 6 ---
\begin{frame}
    \frametitle{Why GMM?}
    
    \vspace{1cm}
    \begin{itemize}
        \item \textbf{PCA alone is insufficient:} If wages affect labor demand, running PCA directly on raw outcomes mixes the endogenous response with the exogenous macro shock.
        \item \textbf{The GMM Solution:} It acts as an iterative filter.
        \begin{itemize}
            \item Guesses the wage effect and subtracts it.
            \item Runs PCA to find the clean macro shock.
            \item Builds the optimal instrument.
            \item Checks if the instrument gives the same wage effect back.
        \end{itemize}
        \item It loops until the system perfectly balances.
    \end{itemize}
\end{frame}

\end{document}
